\documentclass[11pt]{article}
\usepackage[utf8]{inputenc}
\usepackage[T1]{fontenc}
\usepackage[a4paper,margin=1in]{geometry}
\usepackage{lmodern}
\usepackage{microtype}
\usepackage{amsmath,amssymb,amsthm,mathtools}
\usepackage{hyperref}

\title{MathClass: signatures terminales de classes arithmetiques}
\author{Projet Montmory}
\date{\today}

\theoremstyle{plain}
\newtheorem{theorem}{Theoreme}[section]
\newtheorem{proposition}[theorem]{Proposition}
\newtheorem{lemma}[theorem]{Lemme}
\newtheorem{corollary}[theorem]{Corollaire}
\theoremstyle{definition}
\newtheorem{definition}[theorem]{Definition}
\theoremstyle{remark}
\newtheorem{remark}[theorem]{Remarque}

\newcommand{\N}{\mathbb N}
\newcommand{\Pp}{\mathbb P}
\newcommand{\TV}{\mathrm{TV}}
\newcommand{\1}{\mathbf 1}
\newcommand{\F}{\mathcal F}
\newcommand{\G}{\mathcal G}
\newcommand{\M}{\mathcal M}
\newcommand{\push}{_{\#}}

\begin{document}
\maketitle

\begin{abstract}
Cette note formalise une mathematique de classes arithmetiques vues par des
applications terminales. L'exemple directeur est l'entree terminale des
trajectoires de Collatz sous une borne \(B\), mais les resultats principaux sont
des faits finis de theorie de la mesure: decomposition convexe des unions
disjointes, fibres terminales pures, pseudometrique par distance en variation
totale, projectivite entre bornes, et contraction des distances par
push-forward. Ces resultats sont demontres sans supposer la conjecture de
Collatz, Hardy--Littlewood, ni l'infinite des premiers jumeaux. La partie
arithmetique profonde reste l'existence de limites terminales pour des classes
non triviales, notamment les premiers, jumeaux, cousins et sexy primes.
\end{abstract}

\section{Position}

L'objectif est de transformer une question vague,
\[
\text{``Collatz voit-il certaines classes arithmetiques ?''}
\]
en une question mesurable. Une classe arithmetique \(\F\subseteq\N\) est envoyee
sur une loi empirique terminale
\[
\F\longmapsto \nu_{B,\F,N}.
\]
Deux classes peuvent ensuite etre comparees par
\[
D_{B,N}(\F,\G)=\|\nu_{B,\F,N}-\nu_{B,\G,N}\|_{\TV}.
\]

Les classes visees sont par exemple
\[
\N,\quad 2\N+1,\quad \{n:n\equiv a\pmod q\},\quad \Pp,
\]
et les motifs premiers
\[
\mathcal P_H=\{p:p+h\text{ est premier pour tout }h\in H\}.
\]
Les jumeaux correspondent a \(H=\{0,2\}\), les cousins a \(H=\{0,4\}\), les
sexy primes a \(H=\{0,6\}\).

\section{Observation terminale}

\begin{definition}[Application terminale abstraite]
Soit \(X_B\) un ensemble fini ou denombrable et soit
\[
E_B:\N\to X_B
\]
une application. Pour une classe \(\F\subseteq\N\), on pose
\[
\F_N=\F\cap[1,N].
\]
Si \(\F_N\ne\varnothing\), la signature terminale empirique de \(\F\) est
\[
\nu_{B,\F,N}
=
(E_B)\push\left(
\frac{1}{|\F_N|}
\sum_{n\in\F_N}\delta_n
\right).
\]
\end{definition}

\begin{definition}[Cas Collatz accelere]
On considere
\[
T(n)=
\begin{cases}
n/2,& n\equiv0\pmod2,\\
(3n+1)/2,& n\equiv1\pmod2.
\end{cases}
\]
Pour \(B\ge1\),
\[
\tau_B(n)=\inf\{t\ge0:T^t(n)\le B\}.
\]
Si l'ensemble est vide, \(\tau_B(n)=\infty\). On definit
\[
E_B(n)=
\begin{cases}
T^{\tau_B(n)}(n),& \tau_B(n)<\infty,\\
\infty,& \tau_B(n)=\infty.
\end{cases}
\]
Alors
\[
X_B=\{1,\ldots,B\}\cup\{\infty\}.
\]
\end{definition}

\begin{remark}
L'etat \(\infty\) est important: toutes les definitions restent valides sans
supposer la conjecture globale de Collatz.
\end{remark}

\section{Algebre finie des classes}

\begin{proposition}[Decomposition des unions disjointes]
Soient \(\F,\G\subseteq\N\) disjointes. Si
\((\F\cup\G)_N\ne\varnothing\), alors
\[
\nu_{B,\F\cup\G,N}
=
\alpha_N\nu_{B,\F,N}+(1-\alpha_N)\nu_{B,\G,N},
\]
ou
\[
\alpha_N=\frac{|\F_N|}{|\F_N|+|\G_N|},
\]
avec la convention naturelle lorsque l'une des deux classes est vide.
\end{proposition}

\begin{proof}
La mesure uniforme sur \((\F\cup\G)_N\) est la combinaison convexe des mesures
uniformes sur \(\F_N\) et \(\G_N\), avec les poids indiques. Le push-forward par
\(E_B\) est lineaire sur les mesures finies.
\end{proof}

\begin{corollary}[Partitions finies]
Si \(\F=\bigsqcup_{i=1}^r \F_i\), alors
\[
\nu_{B,\F,N}
=
\sum_{i=1}^r
\frac{|(\F_i)_N|}{|\F_N|}\nu_{B,\F_i,N}.
\]
\end{corollary}

\begin{proposition}[Fibres terminales pures]
Pour \(y\in X_B\), soit
\[
A_y(B)=\{n:E_B(n)=y\}.
\]
Si \(A_y(B)\cap[1,N]\ne\varnothing\), alors
\[
\nu_{B,A_y(B),N}=\delta_y.
\]
\end{proposition}

\begin{proof}
Par definition de \(A_y(B)\), tout element observe dans cette classe a image
terminale \(y\). Le push-forward de la mesure uniforme est donc la masse de
Dirac en \(y\).
\end{proof}

\begin{corollary}[Classes definies par un ensemble terminal]
Si \(A\subseteq X_B\) et \(\mathcal C_A=E_B^{-1}(A)\), alors
\[
\nu_{B,\mathcal C_A,N}(A)=1
\]
des que \((\mathcal C_A)_N\ne\varnothing\).
\end{corollary}

\section{Distances de classes}

\begin{definition}[Distance terminale finie]
Pour deux classes \(\F,\G\) ayant des echantillons non vides a \(N\),
\[
D_{B,N}(\F,\G)
=
\|\nu_{B,\F,N}-\nu_{B,\G,N}\|_{\TV}.
\]
\end{definition}

\begin{proposition}[Pseudometrique]
Pour \(B,N\) fixes, \(D_{B,N}\) est une pseudometrique sur les classes
admissibles:
\[
D_{B,N}(\F,\G)\ge0,
\]
\[
D_{B,N}(\F,\G)=D_{B,N}(\G,\F),
\]
et
\[
D_{B,N}(\F,\mathcal H)\le
D_{B,N}(\F,\G)+D_{B,N}(\G,\mathcal H).
\]
\end{proposition}

\begin{proof}
Ces trois proprietes sont celles de la distance en variation totale sur les
mesures de probabilite. Le terme ``pseudo'' est necessaire: deux classes
distinctes peuvent avoir la meme loi terminale.
\end{proof}

\section{Projectivite entre bornes}

\begin{definition}[Projection terminale]
Pour \(B_1<B_2\), on definit
\[
\delta_{B_1,B_2}:X_{B_2}\to X_{B_1}
\]
par
\[
\delta_{B_1,B_2}(\infty)=\infty
\]
et, pour \(y\le B_2\),
\[
\delta_{B_1,B_2}(y)=E_{B_1}(y).
\]
\end{definition}

\begin{theorem}[Projectivite finie exacte]
Pour tout \(n\in\N\),
\[
E_{B_1}(n)=\delta_{B_1,B_2}(E_{B_2}(n)).
\]
Par consequent, pour toute classe \(\F\) et tout \(N\) admissible,
\[
\nu_{B_1,\F,N}
=
(\delta_{B_1,B_2})\push\nu_{B_2,\F,N}.
\]
\end{theorem}

\begin{proof}
Si \(E_{B_2}(n)=\infty\), l'orbite ne rencontre pas \([1,B_2]\), donc elle ne
rencontre pas \([1,B_1]\). Ainsi \(E_{B_1}(n)=\infty\). Sinon
\(E_{B_2}(n)=y\le B_2\). Pour entrer sous \(B_1\), l'orbite partant de \(n\)
atteint d'abord \(y\), puis poursuit comme l'orbite partant de \(y\). L'entree
sous \(B_1\) est donc \(E_{B_1}(y)\). L'identite des mesures suit par
push-forward.
\end{proof}

\begin{lemma}[Contraction par push-forward]
Pour toute application \(f:X\to Y\) et toutes probabilites \(\mu,\nu\) sur \(X\),
\[
\|f\push\mu-f\push\nu\|_{\TV}\le \|\mu-\nu\|_{\TV}.
\]
\end{lemma}

\begin{proof}
Pour tout \(A\subseteq Y\),
\[
(f\push\mu)(A)-(f\push\nu)(A)=\mu(f^{-1}(A))-\nu(f^{-1}(A)).
\]
Le supremum sur les ensembles \(A\subseteq Y\) est donc borne par le supremum
sur tous les sous-ensembles de \(X\).
\end{proof}

\begin{corollary}[Monotonie des distances]
Pour \(B_1<B_2\),
\[
D_{B_1,N}(\F,\G)\le D_{B_2,N}(\F,\G).
\]
\end{corollary}

\section{Passage conditionnel aux limites}

\begin{definition}[Loi terminale limite]
Si la limite existe, on ecrit
\[
\nu_{B,\F,N}\longrightarrow \nu_{B,\F}.
\]
\end{definition}

\begin{proposition}[Identites limites]
Si les lois limites des classes d'une partition finie existent et si les poids
asymptotiques existent, les decompositions convexes finies passent a la limite.
Si \(\nu_{B_2,\F}\) existe, alors pour tout \(B_1<B_2\),
\[
\nu_{B_1,\F}
=
(\delta_{B_1,B_2})\push\nu_{B_2,\F}.
\]
\end{proposition}

\begin{proof}
La premiere assertion est le passage a la limite dans une somme finie. La
seconde suit de la projectivite finie et de la continuite du push-forward sur
un espace de mesures fini.
\end{proof}

\section{Critere par densites naturelles}

Le critere suivant est le pont le plus direct vers une mathematique de classes.
Il transforme l'existence d'une loi terminale en question de densites
d'intersections.

\begin{theorem}[Critere de densite]
Supposons \(X_B\) fini. Pour \(y\in X_B\), posons
\[
A_y(B)=\{n:E_B(n)=y\}.
\]
Soit \(\F\subseteq\N\) une classe de densite naturelle positive \(d(\F)>0\).
Supposons que chaque intersection \(\F\cap A_y(B)\) possede une densite
naturelle
\[
d_y(\F,B)=d(\F\cap A_y(B)).
\]
Alors la loi terminale limite existe et
\[
\nu_{B,\F}(y)=\frac{d_y(\F,B)}{d(\F)}.
\]
\end{theorem}

\begin{proof}
Pour chaque \(y\),
\[
\nu_{B,\F,N}(y)=
\frac{|\F\cap A_y(B)\cap[1,N]|}{|\F\cap[1,N]|}
=
\frac{|\F\cap A_y(B)\cap[1,N]|/N}{|\F\cap[1,N]|/N}.
\]
Le numerateur converge vers \(d_y(\F,B)\) et le denominateur vers \(d(\F)>0\).
Comme \(X_B\) est fini, la convergence coordonnee donne la convergence de la
mesure.
\end{proof}

\begin{corollary}[Neutralite terminale]
Supposons que \(\nu_{B,\N}\) existe. Si
\[
d(\F\cap A_y(B))=d(\F)d(A_y(B))
\]
pour tout \(y\in X_B\), alors
\[
\nu_{B,\F}=\nu_{B,\N}.
\]
\end{corollary}

\begin{proof}
Par le critere,
\[
\nu_{B,\F}(y)=\frac{d(\F)d(A_y(B))}{d(\F)}=d(A_y(B)).
\]
Mais \(d(A_y(B))=\nu_{B,\N}(y)\).
\end{proof}

\begin{proposition}[Forme biaisee]
Supposons que \(\nu_{B,\N}\) existe et que des poids \(w_y(\F)\ge0\) verifient
\[
d(\F\cap A_y(B))=w_y(\F)d(\F)d(A_y(B))
\]
pour tout \(y\), avec normalisation non nulle. Alors
\[
\nu_{B,\F}(y)=
\frac{w_y(\F)\nu_{B,\N}(y)}
{\sum_z w_z(\F)\nu_{B,\N}(z)}.
\]
\end{proposition}

\begin{proof}
La formule brute donne \(\nu_{B,\F}(y)=w_y(\F)d(A_y(B))\). Comme les masses
doivent sommer a \(1\), on ecrit la meme relation sous forme normalisee.
\end{proof}

\section{Motifs premiers}

\begin{definition}[Classe d'un motif premier]
Pour un motif fini \(H\subset\mathbb Z\), on pose
\[
\mathcal P_H=\{p:p+h\text{ est premier pour tout }h\in H\}.
\]
\end{definition}

\begin{theorem}[Loi terminale conditionnelle pour un motif]
Supposons que, pour chaque \(y\in X_B\),
\[
|\{p\le x:p\in\mathcal P_H,\ E_B(p)=y\}|
\sim
c_y(B,H)\frac{x}{(\log x)^{|H|}},
\]
et que
\[
C_B(H)=\sum_{y\in X_B}c_y(B,H)>0.
\]
Alors
\[
\nu_{B,\mathcal P_H,N}(y)\longrightarrow
\frac{c_y(B,H)}{C_B(H)}.
\]
\end{theorem}

\begin{proof}
Le denominateur est la somme finie des numerateurs sur \(y\in X_B\). En
utilisant les equivalents supposes et \(C_B(H)>0\), le quotient converge vers
\(c_y(B,H)/C_B(H)\).
\end{proof}

\begin{remark}
Ce theoreme ne prouve aucune asymptotique de nombres premiers. Il dit seulement
que si les asymptotiques locales terminales existent, alors la loi terminale est
leur normalisation.
\end{remark}

\section{Filtrage terminal et normalisation}

Soit \(A\subseteq X_B\) et
\[
M_A(n)=1\quad\Longleftrightarrow\quad E_B(n)\in A.
\]
Si \(\nu_{B,\mathcal P_H}\) existe, on pose
\[
\rho_A(H)=\nu_{B,\mathcal P_H}(A).
\]
Sous une asymptotique de Hardy--Littlewood
\[
|\mathcal P_H\cap[1,x]|
\sim
S(H)\frac{x}{(\log x)^{|H|}},
\]
le coefficient conditionnel du filtre terminal est
\[
S(H)\rho_A(H).
\]
Pour les jumeaux, \(S(H)=2C_2\), donc
\[
C_{\rm Montmory}=2C_2\,\nu_{B,\rm twin}(A)
\]
dans la convention ou \(C_{\rm Montmory}\) designe le coefficient devant
\(x/\log^2x\). Dans la convention relative,
\[
C_{\rm Montmory}^{\rm rel}=\nu_{B,\rm twin}(A).
\]

\section{Ce qui est obtenu}

On obtient une mathematique finie, exacte et projective des classes:
\[
\boxed{\F\longmapsto \nu_{B,\F,N}}
\]
et
\[
\boxed{(\F,\G)\longmapsto D_{B,N}(\F,\G).}
\]
Elle s'applique aux premiers, jumeaux, cousins et sexy primes comme classes,
mais ne demontre pas encore leurs lois limites. Le verrou mathematique fort est
donc:
\[
\boxed{\text{prouver l'existence de }\nu_{B,\F}
\text{ pour des classes arithmetiques non triviales}.}
\]

\end{document}
